\chapter{Mohs Micrographic Surgery}

\section*{What Is This Surgery?}
Mohs micrographic surgery (MMS) is a highly specialized, tissue-sparing surgical technique for the precise removal of certain skin cancers, primarily basal cell carcinoma (BCC) and squamous cell carcinoma (SCC). This procedure is unique in that the surgeon acts as both the excisional surgeon and the pathologist, meticulously examining 100\% of the surgical margins of each excised tissue layer under a microscope during the operation. This real-time, margin-controlled approach allows for the complete eradication of cancerous cells while preserving the maximum amount of healthy surrounding tissue. Mohs surgery is particularly indicated for skin cancers located in cosmetically or functionally sensitive areas, such as the face, eyelids, nose, ears, lips, and digits, or for tumors with aggressive features, offering the highest reported cure rates for these challenging lesions. Its clinical significance lies in its ability to achieve superior oncologic outcomes with optimal functional and aesthetic preservation.

\section*{Alternative Names}
\begin{itemize}
  \item Mohs micrographic surgery (MMS)
  \item Mohs chemosurgery (historical term)
  \item Microscopically controlled excision
  \item Micrographic surgery
  \item Fresh-tissue Mohs technique
  \item Margin-controlled skin cancer surgery
  \item Mohs surgery
\end{itemize}

\section*{Relevant Surgical Anatomy}
Mohs micrographic surgery primarily targets lesions within the skin and immediate subcutaneous tissues. The skin is composed of three main layers: the epidermis, dermis, and subcutaneous tissue. The epidermis is the outermost protective layer, while the dermis, beneath it, contains connective tissue, blood vessels, nerves, hair follicles, and sweat glands. The subcutaneous tissue (hypodermis) consists mainly of fat and loose connective tissue, providing insulation and cushioning. Surgical considerations vary significantly depending on the anatomical location. On the face, critical structures include branches of the facial nerve (cranial nerve VII), which control facial expression, and sensory nerves like the trigeminal nerve branches (cranial nerve V). The eyelids, nose, and lips are highly specialized structures where even minimal tissue loss can lead to significant functional impairment (e.g., ectropion, nasal airway obstruction, oral incompetence).

The arterial supply to the head and neck is rich, primarily from the external carotid artery branches, necessitating careful hemostasis. Venous drainage generally parallels the arterial supply. Lymphatic drainage patterns are crucial for staging and prognosis, with regional lymph nodes (e.g., preauricular, submandibular, cervical) serving as potential sites for metastatic spread, particularly for aggressive SCC. Surgical landmarks, such as the nasolabial folds, vermilion border of the lip, and free margins of the eyelids, are vital for planning excisions and reconstructions to maintain aesthetic and functional integrity. On the hands and feet, preservation of tendons, nerves, and joints is paramount. The thin skin over cartilage (e.g., ear, nose) and bone (e.g., scalp, shin) presents unique challenges for wound healing and reconstruction due to limited vascularity and tissue laxity.

\section{Surgical Principle \& Rationale}
The fundamental principle of Mohs surgery is the complete microscopic examination of 100\% of the surgical margins to ensure that no malignant cells remain before the surgical defect is repaired. This differs from conventional excision, where margins are typically assessed postoperatively by a separate pathologist, often leading to re-excision if margins are positive. In Mohs, a thin, saucer-shaped layer of tissue is excised, meticulously oriented, and then processed into frozen sections. These sections are then examined by the Mohs surgeon, who maps any residual tumor to its precise location on the patient's wound. If tumor is identified, an additional, highly targeted excision is performed only in the involved area, minimizing the removal of healthy tissue. This iterative process continues until all margins are clear of cancer. The rationale behind this technique is twofold: to achieve the highest possible cure rates for skin cancer by ensuring complete tumor removal, and to maximize the preservation of healthy tissue, thereby optimizing functional and cosmetic outcomes, especially in anatomically complex or cosmetically sensitive regions.

\section{History \& Surgical Evolution}
Mohs surgery was pioneered in the 1930s by Dr. Frederic E. Mohs, a general surgeon at the University of Wisconsin. His initial technique, termed "chemosurgery," involved applying a zinc chloride paste to the tumor, which fixed the tissue in situ, allowing for sequential excision and microscopic examination of the margins. This method, while effective, was time-consuming and caused significant tissue necrosis. A pivotal evolution occurred in the 1960s and 1970s when Dr. Perry Robins and others refined the technique by transitioning to the "fresh-tissue technique." This involved excising unfixed tissue and immediately processing it into frozen sections for microscopic analysis, eliminating the need for the caustic paste and significantly reducing patient discomfort and operative time. This fresh-tissue approach allowed for immediate reconstruction and became the standard for modern Mohs surgery. Subsequent advancements include improved tissue processing techniques, specialized mapping methods, and the integration of immunohistochemical stains for challenging cases, further enhancing the accuracy and applicability of the procedure.

\section{Clinical Indications}
Mohs micrographic surgery is indicated for a range of skin cancers, particularly those with high-risk features or located in critical anatomical areas.

\begin{itemize}
  \item Skin cancers located on the face, eyelids, nose, lips, ears, scalp, neck, hands, feet, shins, nipples, or genitalia, where tissue preservation is critical for function and cosmesis.
  \item Basal cell carcinoma (BCC) and squamous cell carcinoma (SCC) with ill-defined clinical margins, making complete removal challenging with standard excision.
  \item Recurrent skin cancer after previous treatment, including prior excision, radiation, or topical therapies.
  \item Aggressive histologic subtypes, such as infiltrative, morpheaform, micronodular, or basosquamous BCC, and poorly differentiated or desmoplastic SCC.
  \item Large tumors, generally those greater than 2 cm in diameter, or rapidly growing lesions.
  \item Tumors arising in areas of prior radiation, chronic inflammation, or on immunosuppressed patients, who are at higher risk for aggressive disease and recurrence.
  \item Tumors demonstrating perineural or perivascular invasion on biopsy, indicating a higher risk of local spread and recurrence.
  \item Lesions in patients with genetic syndromes predisposing to multiple skin cancers, such as Basal Cell Nevus Syndrome (Gorlin Syndrome) or Xeroderma Pigmentosum.
\end{itemize}

\section{Clinical Positioning}
Mohs micrographic surgery is considered a primary and definitive treatment for appropriately selected skin cancers, particularly those deemed high-risk or located in cosmetically and functionally sensitive areas. It is often the first-line modality due to its superior cure rates and tissue-sparing advantages compared to conventional surgical excision, radiation therapy, or topical treatments. Its role as a primary treatment is especially pronounced for tumors with aggressive histological features, large size, or those with a high likelihood of subclinical extension. Mohs surgery also serves a critical role as a secondary or salvage procedure. This applies when a prior standard excision has resulted in positive or close margins, or when a tumor has recurred after conventional surgery, radiation, or other non-surgical treatments. In these recurrent or histologically aggressive scenarios, Mohs is frequently the preferred retreatment approach because its precise margin control can effectively track and remove microscopic tumor extensions that might be missed by blind re-excision.

\section{Surgical Technique \& Procedure}
Mohs surgery is typically performed in an outpatient setting under local anesthesia. The procedure involves several distinct stages, combining surgical excision with immediate microscopic examination.

\begin{itemize}
  \item \textbf{Anesthesia \& Preparation:} The surgical site is thoroughly cleaned and draped in a sterile fashion. Local anesthetic, typically lidocaine with epinephrine, is infiltrated around and beneath the tumor to ensure patient comfort and provide vasoconstriction for hemostasis.
  \item \textbf{Debulking:} The visible portion of the tumor is often debulked using a curette or scalpel. This step removes the bulk of the tumor, making the subsequent margin-controlled excision more efficient.
  \item \textbf{First Stage Excision:} A thin, saucer-shaped layer of tissue is excised around and beneath the debulked tumor. The angle of excision is typically 45 degrees, creating a beveled edge that facilitates complete margin assessment. The goal is to remove the tumor with a minimal, yet adequate, margin of clinically healthy tissue.
  \item \textbf{Mapping \& Orientation:} The excised tissue specimen is meticulously oriented and marked with colored dyes (e.g., blue, green, black ink) at specific points corresponding to a detailed map drawn of the surgical defect on the patient. This map precisely correlates the tissue specimen to its original anatomical location, allowing the surgeon to pinpoint any residual tumor.
  \item \textbf{Tissue Processing \& Microscopic Examination:} The oriented tissue is immediately transported to an on-site Mohs laboratory. Here, it is flattened, frozen, and sectioned horizontally (en face) to create slides that display 100\% of the peripheral and deep margins. These slides are stained with hematoxylin and eosin (H\&E) and then microscopically examined by the Mohs surgeon.
  \item \textbf{Identification of Residual Tumor:} If malignant cells are identified on the slides, their exact location is marked on the surgical map.
  \item \textbf{Subsequent Stages:} If residual tumor is present, the surgeon returns to the patient and performs another targeted excision, removing only the tissue from the specific area identified on the map as positive. This iterative process of excision, mapping, processing, and microscopic examination continues until all margins are confirmed to be clear of cancer.
  \item \textbf{Defect Repair:} Once clear margins are achieved, the resulting surgical defect is repaired. Options include primary closure (suturing the edges together), secondary intention healing (allowing the wound to heal naturally), or reconstruction using local flaps (moving adjacent skin) or skin grafts (transplanting skin from another body site). The choice of repair depends on the size, location, and depth of the defect, as well as patient factors.
\end{itemize}

\section*{Preoperative Workup \& Preparation}
Preoperative preparation for Mohs micrographic surgery is generally straightforward, as the procedure is typically performed under local anesthesia in an outpatient setting. A thorough medical history and physical examination are essential to assess the patient's overall health status and identify any comorbidities that might affect the procedure or recovery. Patients should provide a comprehensive list of all medications, including prescription drugs, over-the-counter medications, herbal supplements, and vitamins. Special attention is given to anticoagulant or antiplatelet medications (e.g., warfarin, aspirin, clopidogrel), as these can increase the risk of bleeding. The Mohs surgeon will consult with the prescribing physician to determine if these medications can be safely adjusted or temporarily discontinued, weighing the risk of bleeding against the risk of thromboembolic events.

\begin{itemize}
  \item \textbf{Medication Review:} Detailed discussion of all current medications, especially blood thinners.
  \item \textbf{Allergy Assessment:} Identification of any allergies to local anesthetics, latex, or adhesive materials.
  \item \textbf{NPO Status:} Generally, no fasting is required, as local anesthesia is used. Patients are encouraged to eat a normal breakfast.
  \item \textbf{Hygiene:} Patients are advised to shower and wash their hair the morning of surgery, avoiding makeup, lotions, or creams on the surgical site.
  \item \textbf{Transportation:} Patients should arrange for someone to drive them home, especially if an anxiolytic or mild sedative is administered.
  \item \textbf{Prophylactic Antibiotics:} Rarely indicated, but may be considered for patients with prosthetic heart valves, certain orthopedic implants, or those at high risk of infection, in consultation with guidelines.
  \item \textbf{Informed Consent:} Detailed discussion of the procedure, potential risks, benefits, alternatives, and expected outcomes, including the possibility of multiple stages and various reconstructive options.
\end{itemize}

\section*{Contraindications \& Precautions}
While Mohs micrographic surgery is highly effective and generally safe, certain conditions may contraindicate its use or necessitate specific precautions.

\textbf{Absolute Contraindications:}
\begin{itemize}
  \item Patient refusal or inability to cooperate with the procedure, which requires the patient to remain still for extended periods.
  \item Uncontrolled severe bleeding diathesis that cannot be managed preoperatively.
  \item Anaphylactic allergy to local anesthetics that cannot be mitigated.
  \item Tumors with extensive deep invasion (e.g., into bone, orbit, or brain) where the extent of disease makes complete removal by Mohs technically unfeasible or where systemic therapy is more appropriate.
\end{itemize}

\textbf{Relative Contraindications \& Precautions:}
\begin{itemize}
  \item Patients with severe anxiety or claustrophobia who may not tolerate the prolonged procedure, even with oral anxiolytics.
  \item Lesions in anatomical areas where the tissue defect configuration makes closure extremely difficult or where vital structures are at very high risk of injury (e.g., deep parotid gland involvement).
  \item Patients with significant cardiac comorbidities who may not tolerate the epinephrine in local anesthetics or the stress of a prolonged procedure.
  \item Tumors that are clearly benign or low-risk and can be adequately treated with simpler, less resource-intensive methods.
  \item Patients on high-dose anticoagulation where temporary cessation is deemed too risky by their primary physician or cardiologist.
  \item Immunosuppressed patients require particular attention due to their increased risk of aggressive tumors, higher recurrence rates, and potential for impaired wound healing.
  \item Consideration of the likely size of the final scar and planning for potential staged reconstruction or referral to a plastic surgeon for complex defects.
  \item Tumors with suspected distant metastasis, where the primary focus shifts to systemic therapy rather than local excision.
\end{itemize}

\section*{Complications \& Risks}
As with any surgical procedure, Mohs micrographic surgery carries potential risks and complications, which are generally low but important for patients to understand.

\textbf{General Surgical Complications:}
\begin{itemize}
  \item \textbf{Bleeding and Hematoma:} Minor bleeding is common, but significant bleeding or hematoma formation requiring intervention can occur.
  \item \textbf{Infection:} Although sterile technique is used, wound infection is a risk, typically managed with antibiotics.
  \item \textbf{Pain and Tenderness:} Postoperative pain is usually mild and managed with over-the-counter analgesics.
  \item \textbf{Anesthesia-related Risks:} Reactions to local anesthetics (e.g., allergic reaction, systemic toxicity from overdose, vasovagal syncope) are rare but possible.
\end{itemize}

\textbf{Procedure-Specific Complications:}
\begin{itemize}
  \item \textbf{Scarring and Cosmetic Deformity:} While Mohs aims to minimize tissue removal, a scar will always result. The appearance varies based on location, size of defect, and individual healing.
  \item \textbf{Numbness or Altered Sensation:} Injury to superficial sensory nerves can cause temporary or permanent numbness, tingling, or hypersensitivity in the area.
  \item \textbf{Motor Nerve Injury:} In areas like the face, injury to motor nerves (e.g., facial nerve branches) can lead to temporary or permanent weakness or paralysis of facial muscles, though this is rare with careful technique.
  \item \textbf{Wound Healing Problems:} Delayed healing, wound dehiscence, or poor cosmetic outcome can occur, especially over cartilage, bone, or in poorly vascularized areas.
  \item \textbf{Recurrence:} Despite the high cure rates, a small risk of local tumor recurrence exists, particularly for very aggressive tumors or those with extensive subclinical spread.
  \item \textbf{Keloid or Hypertrophic Scarring:} Some individuals are prone to developing raised, thickened scars.
  \item \textbf{Pigmentary Changes:} The treated area may become lighter (hypopigmentation) or darker (hyperpigmentation) than the surrounding skin.
  \item \textbf{Allergic Reactions:} To suture materials, topical antibiotics, or wound dressings.
  \item \textbf{Cartilage or Bone Exposure/Necrosis:} In deep excisions, especially on the nose or ear, cartilage or bone may be exposed, increasing the risk of infection or necrosis.
\end{itemize}

\section*{Postoperative Recovery Timeline}
Immediately after Mohs surgery, patients are monitored briefly in the recovery area to ensure stable vital signs and adequate pain control. The surgical wound is dressed, and detailed wound care instructions are provided. During the first 24-48 hours, patients typically experience mild pain, swelling, and bruising around the surgical site, which can be managed with over-the-counter pain relievers and cold compresses. Activity restrictions, such as avoiding strenuous exercise, heavy lifting, and bending, are usually advised to minimize bleeding and swelling. The dressing typically needs to be kept clean and dry, with specific instructions for dressing changes provided.

Over the first 1-2 weeks, the wound continues to heal. Sutures, if used, are generally removed within 7-14 days, depending on the anatomical location and tension on the wound. Patients are advised to continue gentle wound care, avoid direct sun exposure, and protect the healing site. Swelling and bruising gradually subside, and the initial redness of the scar begins to fade. Full scar maturation is a prolonged process, taking several months to a year or more, during which the scar will soften, flatten, and lighten in color. Patients are encouraged to protect the healing area from trauma and sun exposure to optimize cosmetic outcomes and reduce the risk of future skin cancers.

\section*{Surgical Findings \& Pathology}
During Mohs micrographic surgery, the surgeon meticulously documents several key findings. This includes the initial clinical size and appearance of the tumor, its anatomical location, and the number of stages required to achieve clear margins. The precise mapping of each excised tissue layer and the corresponding areas of residual tumor are critical surgical findings. The pathology report, generated from the microscopic examination of the resected tissues, provides definitive confirmation of the tumor type (e.g., basal cell carcinoma, squamous cell carcinoma), its histological subtype (e.g., nodular, infiltrative, morpheaform), and any adverse features such as perineural or perivascular invasion. Crucially, the pathology report confirms the status of the surgical margins for each stage. A "clear margin" on the final stage signifies that no cancerous cells were detected at the edges of the excised tissue, indicating complete tumor removal. The combination of surgical findings and the detailed pathology report guides subsequent patient management, including the choice of reconstruction and the need for further surveillance or adjuvant therapies.

\section*{Expected Postoperative Course}
A normal and successful postoperative course following Mohs micrographic surgery is characterized by progressive wound healing without significant complications. Patients can expect mild to moderate pain and discomfort in the first few days, which should gradually resolve with conservative management. Swelling and bruising are common and typically subside within one to two weeks. The surgical wound will gradually close and mature, with the initial redness and firmness of the scar diminishing over several months. The goal is for the wound to heal into an aesthetically acceptable scar that blends with the surrounding skin, minimizing functional impairment. There should be no signs of infection, such as increasing redness, warmth, pus, or fever. Ultimately, a successful outcome means complete eradication of the skin cancer, confirmed by clear margins on the final pathology, and no evidence of local recurrence during long-term follow-up. The preserved tissue and precise reconstruction aim to restore both the function and appearance of the treated area.

\section*{Postoperative Care \& Follow-up}
Postoperative care after Mohs surgery focuses on wound management, monitoring for complications, and long-term surveillance for new skin cancers.

\begin{itemize}
  \item \textbf{Wound Care Instructions:} Patients receive detailed instructions on how to clean the wound, change dressings, and protect the site from trauma and sun exposure.
  \item \textbf{Suture Removal:} If sutures were used for closure, they are typically removed by the surgeon or a nurse one to two weeks after surgery, depending on the location and healing progress.
  \item \textbf{Follow-up Visits:} Initial follow-up visits are scheduled to assess wound healing, remove sutures, and address any patient concerns.
  \item \textbf{Scar Management:} Patients may be advised on scar massage, silicone sheeting, or other topical treatments to optimize scar appearance. Physical therapy or occupational therapy is rarely needed unless a motor nerve injury occurred or a large functional defect was created.
  \item \textbf{Sun Protection:} Emphasizing rigorous photoprotection, including sunscreen use, protective clothing, and avoidance of peak sun hours, is crucial to prevent new skin cancers.
  \item \textbf{Regular Skin Examinations:} Patients treated for skin cancer are at increased risk for developing new primary lesions. Therefore, regular full-skin examinations by a dermatologist are recommended, typically every 6-12 months.
  \item \textbf{Warning Signs:} Patients are educated on warning signs of complications (e.g., increasing pain, redness, swelling, pus, fever) or potential recurrence (e.g., new nodule, persistent non-healing sore, change in sensation) and instructed to contact their doctor immediately if these occur.
\end{itemize}
Ongoing follow-up is essential not only to monitor the surgical site but also to screen for new primary skin cancers, which are common in this patient population.

\section*{Epidemiology}
Skin cancer is the most prevalent form of cancer globally, with basal cell carcinoma (BCC) and squamous cell carcinoma (SCC), collectively known as nonmelanoma skin cancers (NMSC), accounting for the vast majority of cases. The incidence of NMSC has been steadily rising over recent decades, particularly in fair-skinned populations. These cancers are strongly associated with cumulative ultraviolet (UV) radiation exposure, fair skin type, advancing age, and male sex. BCC is the most common type of skin cancer, typically presenting on sun-exposed areas of the head and neck, which are frequent sites for Mohs surgery. SCC is the second most common, with a higher potential for metastasis, especially in high-risk locations or in immunosuppressed individuals. Immunosuppression, prior radiation therapy, and certain genetic syndromes (e.g., Gorlin syndrome, xeroderma pigmentosum) significantly increase the risk of developing aggressive NMSC. Mohs micrographic surgery has become one of the most frequently performed procedures for high-risk NMSC due to its efficacy and tissue-sparing benefits, reflecting the increasing burden of these cancers. While mortality from NMSC is low overall, morbidity can be significant due to local tissue destruction and potential for metastasis in aggressive cases.

\section*{Outcomes \& Prognosis}
Mohs micrographic surgery offers the highest reported cure rates for basal cell carcinoma and squamous cell carcinoma, making it the gold standard for high-risk lesions. Five-year cure rates typically exceed 95\% for primary BCC and 99\% for primary SCC in many large series, significantly outperforming conventional excision for these indications. For recurrent tumors, while cure rates are slightly lower, Mohs surgery still demonstrates superior efficacy compared to other modalities. The primary functional and cosmetic outcomes are generally excellent due to the maximal preservation of healthy tissue, which is particularly critical for lesions on the face, hands, and genitalia. Prognostic factors influencing recurrence risk include tumor size, location, histological subtype (e.g., infiltrative, perineural invasion), and the patient's immune status. For instance, SCC with perineural invasion or in immunosuppressed patients carries a higher risk of recurrence and metastasis, necessitating closer follow-up. Landmark studies, such as those comparing Mohs to standard excision for high-risk NMSC, consistently demonstrate Mohs' superior efficacy in terms of local control. Overall, the prognosis for patients treated with Mohs surgery is excellent, provided the tumor is completely removed and new lesions are detected early through diligent surveillance.

\section*{References}
\begin{enumerate}
  \item American Academy of Dermatology. Mohs Micrographic Surgery Guidelines of Care. J Am Acad Dermatol. 2024;90(1):1-20.
  \item Schwartz's Principles of Surgery. 12th ed. McGraw-Hill Education; 2023.
  \item Mohs FE. Chemosurgery: a microscopically controlled method of cancer excision. Arch Surg. 1941;42(2):279-295.
  \item MedlinePlus. Mohs Surgery. U.S. National Library of Medicine; 2025.
\end{enumerate}

\clearpage